\documentclass[11pt,a4paper]{article}

% ---------------------------------------------------------------
%  Teaching Statement -- Stanisław Szufa
%  EPFL IC Faculty Search 2026
% ---------------------------------------------------------------
\usepackage[T1]{fontenc}
\usepackage[utf8]{inputenc}
\usepackage{mathpazo}
\renewcommand{\ttdefault}{lmtt}  % outline typewriter font (avoids bitmap ectt)
\usepackage[margin=2.5cm]{geometry}
\usepackage{microtype}
\usepackage{enumitem}
\usepackage{titlesec}
\usepackage{xcolor}
\usepackage[hidelinks]{hyperref}

\definecolor{accent}{RGB}{160,20,30}   % same accent as the research statement
\titleformat{\section}{\large\bfseries\color{accent}}{}{0pt}{}
\titlespacing*{\section}{0pt}{9pt}{3pt}
\titleformat{\subsection}[runin]{\normalsize\bfseries}{}{0pt}{}[.]
\titlespacing*{\subsection}{0pt}{5pt}{0.6em}
\setlength{\parindent}{0pt}
\setlength{\parskip}{4pt}
\setlist{nosep,leftmargin=1.2em}

\begin{document}

\begin{center}
  {\LARGE\bfseries Teaching Statement}\\[3pt]
  {\large Stanisław Szufa}\\[1pt]
  {\small University of Geneva}
\end{center}
\vspace{2pt}

% ===============================================================
\section{Teaching Philosophy: Learning by Measuring}
% ===============================================================

The most useful thing I learned during my PhD, I had to learn on my own.
My advisor's work was theoretical, and when I wanted to know how voting rules
actually behave, nobody around me could tell me how to design a numerical
experiment, which data to use, or how to tell a real finding from an artifact
of the generator. Working this out became my thesis, and then a tutorial, a
community guide, and two software libraries, so that others would not have to
start from scratch. This experience shapes how I want to teach. I want students
to \emph{measure} the objects they study, not only to prove things about them:
a student who has computed an optimal consensus ranking for a real election
and watched an integer-programming solver stall on it understands NP-hardness
in a way that no reduction on a blackboard can give. The evidence on active
learning in STEM points the same way (Freeman et al., \emph{PNAS}, 2014).

When I design my own courses, this will mean three things:
\begin{itemize}
  \item \emph{Every abstraction comes with an instance.} I want to introduce
  each definition together with a dataset on which students can try it in the
  same session. In my field this is easy: a city's participatory budget is a
  real election with real stakes, and students can load one from Pabulib, run
  three voting rules on it, and argue about which outcome is fairer before
  they have seen a single axiom.
  \item \emph{Students should discover, not only verify.} Standard exercises
  build the foundations, but the course projects I want to set have answers
  that I do not know in advance. ``Which random model of voters best
  approximates the preferences of Warsaw's residents?'' is a question a
  second-year student can attack with the tools from my research, and, as the
  student projects I have co-supervised show, the answers can be publishable.
  \item \emph{Tools are part of the curriculum.} Writing clean, reproducible
  experimental code is a skill that most students are never taught. I intend
  to teach it explicitly, with open-source libraries such as
  \texttt{prefsampling} and \texttt{mapof} as real codebases that students
  read, test, and extend instead of writing throwaway scripts.
\end{itemize}

% ===============================================================
\section{Teaching Experience}
% ===============================================================

As a Teaching and Research Assistant at Jagiellonian University (2018--2021),
I taught the exercise sessions of three core bachelor courses in computer
science: \emph{Introduction to Programming}, \emph{Logic and Set Theory}, and
\emph{Linear Algebra}. These were courses designed by others, with large
first- and second-year cohorts and a wide spread of preparation, and they
taught me to explain the same idea in several ways and to notice early which
students are falling behind. % TODO: add cohort sizes, number of groups/sessions, evaluation scores or student feedback if available
They also cover the material I would most like to teach at the bachelor
level at EPFL: programming, discrete mathematics, and algorithms are the
foundation on which everything in my research stands.

My own research I have so far taught only outside the classroom, to my
research community. I gave a tutorial on numerical experiments in
computational social choice at three of the major AI conferences, AAAI,
AAMAS, and IJCAI, % TODO: add years
and I co-authored the \emph{Guide to Numerical Experiments on Elections}
(IJCAI-24), written as teaching material: it explains what goes wrong in
experiments and how to avoid it.
Other groups now use \texttt{prefsampling} and \texttt{mapof} in their
research and in student projects. Explaining voting rules to city officials
and residents through \emph{Science for Democracy} has been the most
demanding teaching I have done: the audience has no patience for notation and
every question is about a real decision.

% ===============================================================
\section{Mentoring and Supervision}
% ===============================================================

I have co-supervised more than ten bachelor's and master's students, mostly at
AGH University of Kraków, within an ERC-funded project, and at LAMSADE in
Paris. Several of these projects grew into publications at AAAI, AAMAS, ECAI,
and NeurIPS with the student as a coauthor. % TODO: cite 2--3 of these papers (e.g. ECAI'25, AAAI'25, NeurIPS'25) and say where the students went next; per the MIT CommLab guide, avoid naming the students themselves
My approach follows how I was mentored, and how I wish I had been:
\begin{itemize}
  \item \emph{Start small and concrete.} A student's first project is a
  well-defined experimental question with a dataset and a tool already in
  hand, so that the first result comes within weeks.
  \item \emph{Then hand over an open problem.} Once a student has shipped
  something, I give them a question whose answer I do not know, and I expect
  them to disagree with me. The best paper award I received at AAMAS-23 was
  for a paper written by three PhD students without a senior coauthor, and I
  want my own students to have that experience.
  \item \emph{Teach the craft, not only the result.} I review students' code
  and experimental design as carefully as their proofs, and I ask them to
  present their work early and often: clear writing and presenting decide
  whether good work is noticed.
\end{itemize}
As President of the PhD Students' Association at Jagiellonian University and
as a representative in the Polish national association of doctoral
candidates, I spent two years working on funding, supervision rules, and
representation of doctoral students. I learned how much a student's success
depends on things outside the research itself, and I would attend to them in
a group of my own.

% ===============================================================
\section{Teaching at EPFL}
% ===============================================================

\textbf{Existing courses.} I am prepared to teach core courses of the
computer science bachelor, in particular introductory programming, discrete
mathematics (as in \emph{Advanced Information, Computation, Communication~I},
CS-101), \emph{Algorithms~I} (CS-250), and \emph{Theory of Computation}
(CS-251). At the master level, I am prepared to teach game theory, including
\emph{Algorithmic Game Theory} (MGT-435), from equilibria and mechanism
design to the social-choice topics closest to my research.

\textbf{New courses.} I would like to develop two courses that complement
\emph{Technologies for Democratic Society} (CS-234), which studies the systems
behind democratic processes, and \emph{Algorithmic Game Theory}, which
studies equilibria and mechanism design:
\begin{itemize}
  \item \emph{Algorithms for Collective Decisions} (master level). Voting,
  committee elections, participatory budgeting, and fair division, treated as
  algorithmic and data problems: every rule comes with real data and every
  axiom with an experiment. By the end, students should be able to choose an
  aggregation rule for a concrete decision, compute it efficiently, test its
  properties on real and synthetic data, and explain the choice to the people
  it affects. The project is a full study of a real participatory process,
  from data to recommendation.
  \item \emph{AI and Democracy} (master seminar). How language models change
  collective decision-making: AI-mediated deliberation, LLM proxies,
  representation in generated text, and the risks of each. The seminar would
  combine recent papers with hands-on simulation, follow on naturally from
  CS-234, and bring together students of computer science, data science, and
  digital humanities.
\end{itemize}

\textbf{Growing as a teacher.} As my classroom experience is limited to
exercise sessions, I would make full use of the induction seminar and
pedagogical consultations offered by EPFL's Teaching Support Center (CAPE),
and in every course I would ask students for feedback mid-semester, while
there is still time to act on it.

\textbf{Beyond the classroom.} I would bring my work with Science for
Democracy to EPFL through student projects with Swiss municipalities that run
participatory budgets. Teaching residents why a proportional rule funded
one project and not another is, to me, the same job as teaching students why
an algorithm is correct.

\end{document}
